\expandafter\gdef\csname SNside1\endcsname{\node[anchor=north west,inner sep=0pt,text width=48mm,font=\SideFont,align=left] at ([xshift=138mm,yshift=-500.62pt]current page.north west) {\textcolor{SNink}{\textbf{错项排除 · 字符单位不等于字节随机寻址}}\par\vspace{1.5mm}终端按字符流收发，并不意味着可按任意字节地址跳到某个字符；磁盘作为块设备可按块号定位。数据传送单位、是否可寻址、是否独占属于不同分类维度，不能把``字符''与``内存字节地址''机械对应。};}
\expandafter\gdef\csname SNside2\endcsname{\node[anchor=north west,inner sep=0pt,text width=48mm,font=\SideFont,align=left] at ([xshift=138mm,yshift=-397.86pt]current page.north west) {\textcolor{SNink}{\textbf{课程连接 · DMA与中断可以协作}}\par\vspace{1.5mm}组成原理说明DMA如何取得存储访问机会并搬运数据，操作系统说明驱动如何设置地址长度以及怎样处理完成。DMA结束常用中断通知，二者分别承担传送与通知，不能因为DMA无需逐字搬运就排除中断或内核参与。};}
\expandafter\gdef\csname SNside3\endcsname{\node[anchor=north west,inner sep=0pt,text width=48mm,font=\SideFont,align=left] at ([xshift=138mm,yshift=-95.00pt]current page.north west) {\textcolor{SNink}{\textbf{错项排除 · 驱动知道设备细节}}\par\vspace{1.5mm}文件逻辑位置属于上层抽象，控制器寄存器、传统CHS参数和启动命令属于设备细节。驱动在提交时形成具体请求，中断处理在完成或异常时推进状态。若题设改为设备内部转换LBA，应按新接口重新分配职责。};}
\expandafter\gdef\csname SNside4\endcsname{\node[anchor=north west,inner sep=0pt,text width=48mm,font=\SideFont,align=left] at ([xshift=138mm,yshift=-95.00pt]current page.north west) {\textcolor{SNink}{\textbf{最小反例 · 不阻塞仍可能忙等}}\par\vspace{1.5mm}程序不断调用非阻塞read，得到暂不可用后立即再次调用，虽然每次调用都不睡眠，CPU却可能一直空转。接口不阻塞只描述本次调用的等待行为，不能直接推出自动提高利用率；异步完成通知另有接口语义。};}
\expandafter\gdef\csname SNside5\endcsname{\node[anchor=north west,inner sep=0pt,text width=48mm,font=\SideFont,align=left] at ([xshift=138mm,yshift=-95.00pt]current page.north west) {\textcolor{SNink}{\textbf{不变量 · 缓冲复用不能覆盖未用数据}}\par\vspace{1.5mm}设备只有在对应系统缓冲已可复用时才能填下一块，复制也不能覆盖仍被分析的用户工作区。按这两个约束检查每段重叠，再求max形式的周期；若复制不用CPU或存在多个工作区，原公式的资源模型已改变。};}
\expandafter\gdef\csname SNside6\endcsname{\node[anchor=north west,inner sep=0pt,text width=48mm,font=\SideFont,align=left] at ([xshift=138mm,yshift=-95.00pt]current page.north west) {\textcolor{SNink}{\textbf{最小反例 · 一块数据没有流水收益}}\par\vspace{1.5mm}只有一个块时，输入T、复制M、分析C仍须完成，单缓冲与双缓冲同为T＋M＋C。双缓冲的收益来自后续块输入与前一块CPU工作重叠，不能把``增加一个缓冲区''直接解释成任何工作量都耗时减半。};}
\expandafter\gdef\csname SNside7\endcsname{\node[anchor=north west,inner sep=0pt,text width=48mm,font=\SideFont,align=left] at ([xshift=138mm,yshift=-95.00pt]current page.north west) {\textcolor{SNink}{\textbf{课程连接 · 设备链也可能形成资源环}}\par\vspace{1.5mm}进程拿到设备后再等控制器，另一进程持控制器又等所需设备，就可能形成占有并等待与循环等待。设备分配和进程资源分配共享死锁分析方法；仅把某张设备表标为已分配，并不代表完整传输路径已经可用。};}
\expandafter\gdef\csname SNside8\endcsname{\node[anchor=north west,inner sep=0pt,text width=48mm,font=\SideFont,align=left] at ([xshift=138mm,yshift=-95.00pt]current page.north west) {\textcolor{SNink}{\textbf{最小反例 · 最高请求不一定就是端点}}\par\vspace{1.5mm}磁头58向上，剩余请求含180和15，物理端点199。LOOK可在180折返，SCAN按端点规则继续到199再折返，二者距离不同。原配套题恰好请求199只能说明那组数据相同，不能推出两算法普遍等价。\par\smallskip\begin{tikzpicture}[x=1mm,y=1mm,>=Stealth,node distance=5mm]
\node[draw=SNink!50,fill=SNblue,rounded corners=1mm,text width=34mm,align=center,inner sep=1.2mm,font=\fontsize{7.5}{9}\selectfont] (m0) at (0,0) {58向上};
\node[draw=SNink!50,fill=SNblue,rounded corners=1mm,text width=34mm,align=center,inner sep=1.2mm,font=\fontsize{7.5}{9}\selectfont] (m1) at (0,-11) {最高180};
\node[draw=SNink!50,fill=SNblue,rounded corners=1mm,text width=34mm,align=center,inner sep=1.2mm,font=\fontsize{7.5}{9}\selectfont] (m2) at (0,-22) {端点199};
\draw[->,SNink] (m0) -- (m1);
\draw[->,SNink] (m1) -- (m2);
\end{tikzpicture}
\par{\fontsize{7.5}{9}\selectfont 仍有15号请求时，SCAN到199折返，LOOK可在180折返。}};}
\expandafter\gdef\csname SNside9\endcsname{\node[anchor=north west,inner sep=0pt,text width=48mm,font=\SideFont,align=left] at ([xshift=138mm,yshift=-95.00pt]current page.north west) {\textcolor{SNink}{\textbf{错项排除 · SSTF每走一步重算距离}}\par\vspace{1.5mm}初始磁头85，服务100以后，剩余101只差1，而60差40。算法比较的是距当前磁头的距离，不能一次按距初始85的差值排完。有限固定批次最终会完成，``持续到达可能饥饿''也不等于任何批次必然饥饿。};}
\expandafter\gdef\csname SNside10\endcsname{\node[anchor=north west,inner sep=0pt,text width=48mm,font=\SideFont,align=left] at ([xshift=138mm,yshift=-95.00pt]current page.north west) {\textcolor{SNink}{\textbf{课程连接 · 文件删除不等于立即擦除}}\par\vspace{1.5mm}文件系统取消名字与块使用关系，TRIM通知设备哪些逻辑数据不再需要，FTL再按内部有效页和擦除块组织回收。三个层次的状态变化可以有时间差；闪存通常不能像内存字节那样任意原地改写，也不按单页直接擦除。};}
\expandafter\gdef\csname SNside11\endcsname{\node[anchor=north west,inner sep=0pt,text width=48mm,font=\SideFont,align=left] at ([xshift=138mm,yshift=-95.00pt]current page.north west) {\textcolor{SNink}{\textbf{错项排除 · 逻辑共享没有增加打印头}}\par\vspace{1.5mm}多个作业可同时提交到输出井，由后台服务逐份交给物理打印机。应用看到可共享的服务，但真实打印装置仍一次处理一个任务。SPOOLing把提交与慢设备使用解耦，不能把这种虚拟化说成硬件同时打印多份纸张。};}
\expandafter\gdef\csname SNside12\endcsname{\node[anchor=north west,inner sep=0pt,text width=48mm,font=\SideFont,align=left] at ([xshift=138mm,yshift=-95.00pt]current page.north west) {\textcolor{SNink}{\textbf{不变量 · 旋转延迟先看转速单位}}\par\vspace{1.5mm}7200转每分先除60得到每秒120转，一圈约8.333毫秒，随机角位置平均等待半圈。若题目直接给秒每转，就改用其一半。先把单位写在每个量后面，能避免把r除2与1除2r这两种形式混用。};}
\expandafter\gdef\csname SNside13\endcsname{\node[anchor=north west,inner sep=0pt,text width=48mm,font=\SideFont,align=left] at ([xshift=138mm,yshift=-95.00pt]current page.north west) {\textcolor{SNink}{\textbf{课程连接 · 文件删除不等于立即擦除}}\par\vspace{1.5mm}文件系统取消名字与块使用关系，TRIM通知设备哪些逻辑数据不再需要，FTL再按内部有效页和擦除块组织回收。三个层次的状态变化可以有时间差；闪存通常不能像内存字节那样任意原地改写，也不按单页直接擦除。};}
\expandafter\gdef\csname SNside14\endcsname{\node[anchor=north west,inner sep=0pt,text width=48mm,font=\SideFont,align=left] at ([xshift=138mm,yshift=-95.00pt]current page.north west) {\textcolor{SNink}{\textbf{最小反例 · 最高请求不一定就是端点}}\par\vspace{1.5mm}磁头58向上，剩余请求含180和15，物理端点199。LOOK可在180折返，SCAN按端点规则继续到199再折返，二者距离不同。原配套题恰好请求199只能说明那组数据相同，不能推出两算法普遍等价。\par\smallskip\begin{tikzpicture}[x=1mm,y=1mm,>=Stealth,node distance=5mm]
\node[draw=SNink!50,fill=SNblue,rounded corners=1mm,text width=34mm,align=center,inner sep=1.2mm,font=\fontsize{7.5}{9}\selectfont] (m0) at (0,0) {58向上};
\node[draw=SNink!50,fill=SNblue,rounded corners=1mm,text width=34mm,align=center,inner sep=1.2mm,font=\fontsize{7.5}{9}\selectfont] (m1) at (0,-11) {最高180};
\node[draw=SNink!50,fill=SNblue,rounded corners=1mm,text width=34mm,align=center,inner sep=1.2mm,font=\fontsize{7.5}{9}\selectfont] (m2) at (0,-22) {端点199};
\draw[->,SNink] (m0) -- (m1);
\draw[->,SNink] (m1) -- (m2);
\end{tikzpicture}
\par{\fontsize{7.5}{9}\selectfont 仍有15号请求时，SCAN到199折返，LOOK可在180折返。}};}
\AddToShipoutPictureFG{\begin{tikzpicture}[remember picture,overlay]\ifcsname SNside\thepage\endcsname\csname SNside\thepage\endcsname\fi\end{tikzpicture}}